\documentclass[11pt]{article}

\usepackage[margin=1in]{geometry}
\usepackage{amsmath}
\usepackage{amssymb}
\usepackage{hyperref}

\title{The Current \texttt{Dynamics} Class}
\author{}
\date{}
\setlength\parindent{0cm}

\newcommand{\parnum}{(\arabic{parcount})}

\newcounter{parcount}
\newcommand\p{%
    \stepcounter{parcount}%
    \parnum \hspace{1em}%
}

\newenvironment{parnumbers}{%
   \par%
   \everypar{\noindent \stepcounter{parcount}\parnum \hspace{1em}}%
}{}
\begin{document}

\maketitle

\section{Purpose}

The current \texttt{Dynamics} class is the collision-response kernel for the
clean rebuild. It is intentionally separated from \texttt{Base}. The
\texttt{Base} class owns the simulation loop, particle storage, contact
detection, double-buffered position movement, and reporting. The
\texttt{Dynamics} class owns only the math that converts detected contacts into
momentum changes.

This separation is important because the dynamics code is the part most likely
to be moved later into C++ or GLSL. The current design is therefore written with
a simple kernel-like rule:

\[
    \text{one worker owns one particle}.
\]

During collision processing, the worker for particle \(i\) may read selected
state from another particle \(j\), but it writes only to particle \(i\).

\section{Particle Data Contract}

Each particle is currently stored as a Python dictionary. The dynamics class
expects the following fields to exist before collision processing begins.

\subsection{Position}

Position is double-buffered:

\[
\texttt{location}
=
\{
\texttt{use},\texttt{x1},\texttt{y1},\texttt{x2},\texttt{y2}
\}.
\]

If
\[
    \texttt{use}=1,
\]
then the active position is
\[
    \mathbf{x}=(x_1,y_1).
\]
If
\[
    \texttt{use}=2,
\]
then the active position is
\[
    \mathbf{x}=(x_2,y_2).
\]

The helper \texttt{current\_location()} reads this active position.

\subsection{Particle Properties}

The dynamics code reads:

\[
    r_i = \texttt{radius},
    \qquad
    m_i = \texttt{mass},
    \qquad
    \mathbf{v}_i = (v_{x,i},v_{y,i}).
\]

The current overlap-area formula assumes equal-radius disks:

\[
    r_i = r_j = r.
\]

If unequal radii are passed to \texttt{circle\_overlap\_area()}, the code raises
an error.

\subsection{Contact List}

\texttt{Base.detect\_contacts()} fills each particle's contact storage before
\texttt{Dynamics.process\_collisions()} runs. For particle \(i\), the relevant
fields are:

\[
    \texttt{contact\_count}_i,
    \qquad
    \texttt{contacts}_i.
\]

\texttt{contacts} is a fixed-length array of particle indices. Only the first
\texttt{contact\_count} entries are valid.

\section{Memory Ownership Rule}

The dynamics pass is particle-owned. When processing particle \(i\), the code
loops through \(i\)'s own contacts:

\[
    j \in \texttt{contacts}_i.
\]

The code may read particle \(j\)'s current location and radius to calculate
geometry. It does not write to particle \(j\). This is different from a
traditional pair solver that updates both particles in one pair operation.

The current update has the form:

\[
    \Delta \mathbf{p}_i
    =
    \sum_{j \in C_i}
    \Delta \mathbf{p}_{ij}.
\]

The corresponding update for particle \(j\) is produced only when particle
\(j\)'s own worker processes its own contact list.

\section{Contact Geometry}

For a pair \(i,j\), the active center positions are:

\[
    \mathbf{x}_i = (x_i,y_i),
    \qquad
    \mathbf{x}_j = (x_j,y_j).
\]

The separation vector from particle \(i\) to particle \(j\) is:

\[
    \mathbf{d}_{ij}
    =
    \mathbf{x}_j-\mathbf{x}_i.
\]

The center distance is:

\[
    d = \|\mathbf{d}_{ij}\|.
\]

Contact exists when:

\[
    d < r_i+r_j.
\]

If there is no overlap, \texttt{particle\_contact()} returns
\texttt{None}.

\subsection{Contact Normal}

When \(d>0\), the unit normal from particle \(i\) toward particle \(j\) is:

\[
    \mathbf{n}_{ij}
    =
    \frac{\mathbf{x}_j-\mathbf{x}_i}{d}
    =
    (n_x,n_y).
\]

If the centers are effectively identical,
\[
    d \le 10^{-12},
\]
the code uses the deterministic fallback:

\[
    \mathbf{n}_{ij}=(1,0).
\]

This avoids division by zero.

\section{Penetration Limit Used For Area}

The raw overlap depth is:

\[
    \delta_{\text{raw}}
    =
    r_i+r_j-d.
\]

The current code caps the depth used for area calculation:

\[
    \delta
    =
    \min\left(
        r_i+r_j-d,
        \min(r_i,r_j)
    \right).
\]

Under the equal-radius assumption, this means:

\[
    \delta \le r.
\]

This matches the current development assumption that particles do not move past
the case where the leading edge reaches the other particle center.

The distance sent to the overlap-area formula is:

\[
    d_A
    =
    r_i+r_j-\delta.
\]

\section{Equal-Radius Overlap Area}

For equal-radius particles,
\[
    r_i=r_j=r,
\]
the circular lens area is:

\[
    A(d_A)
    =
    2r^2\cos^{-1}\left(\frac{d_A}{2r}\right)
    -
    \frac{d_A}{2}\sqrt{4r^2-d_A^2}.
\]

The implementation clamps the distance into the valid range:

\[
    d_c
    =
    \min(\max(d_A,0),2r).
\]

The actual computed area is:

\[
    A(d_c)
    =
    2r^2\cos^{-1}\left(\frac{d_c}{2r}\right)
    -
    \frac{d_c}{2}\sqrt{\max(0,4r^2-d_c^2)}.
\]

The \(\max(0,\cdot)\) term protects the square root from small floating-point
roundoff errors.

\section{Overlap Momentum}

The current model follows the minimal legacy Mom-style idea:

\[
    \text{overlap area}
    \rightarrow
    \text{momentum amount}.
\]

It does not treat overlap as a force integrated over the substep. There is no
multiplication by \(\Delta t\) inside \texttt{overlap\_momentum()}.

For a contact between particle \(i\) and particle \(j\), the scalar momentum
amount is:

\[
    J_{ij}
    =
    k_i A_{ij} w(d).
\]

Here:

\[
    k_i = \texttt{momentum\_per\_area}
\]

and:

\[
    A_{ij} = \text{overlap area}.
\]

The coefficient \(k_i\) is owned by particle \(i\). If particle \(i\) has a
\texttt{momentum\_per\_area} field, that value is used. Otherwise the class
default is used. The code also keeps the legacy fallback name
\texttt{repulsion\_force\_per\_area}.

\section{Inverse-Square Weight}

The distance weight is:

\[
    w(d)
    =
    \frac{1}{\max(d^2,s^2)}.
\]

where

\[
    s = \texttt{inverse\_square\_softening}.
\]

The softening value prevents singular behavior when two centers are very close:

\[
    s \ge 10^{-12}.
\]

\section{Momentum Direction}

The contact normal \(\mathbf{n}_{ij}\) points from particle \(i\) toward
particle \(j\). The collision response for particle \(i\) points away from
particle \(j\), so the update for this contact is:

\[
    \Delta \mathbf{p}_{ij}
    =
    -J_{ij}\mathbf{n}_{ij}.
\]

Component-wise:

\[
    \Delta p_{x,ij} = -n_x J_{ij},
    \qquad
    \Delta p_{y,ij} = -n_y J_{ij}.
\]

\section{Accumulating Contacts}

For a particle with contact set \(C_i\), the dynamics class accumulates:

\[
    P_{x,i}
    =
    \sum_{j\in C_i}
    -n_{x,ij}J_{ij},
\]

\[
    P_{y,i}
    =
    \sum_{j\in C_i}
    -n_{y,ij}J_{ij}.
\]

It also tracks the scalar sum:

\[
    P_{\text{sum},i}
    =
    \sum_{j\in C_i}
    J_{ij}.
\]

These values are written back onto particle \(i\) as:

\[
\begin{aligned}
    \texttt{overlap\_momentum\_x} &= P_{x,i},\\
    \texttt{overlap\_momentum\_y} &= P_{y,i},\\
    \texttt{overlap\_momentum} &= P_{\text{sum},i}.
\end{aligned}
\]

The same vector is currently also stored as internal momentum:

\[
\begin{aligned}
    \texttt{internal\_momentum\_x} &= P_{x,i},\\
    \texttt{internal\_momentum\_y} &= P_{y,i},\\
    \texttt{internal\_momentum} &=
    \sqrt{P_{x,i}^2+P_{y,i}^2}.
\end{aligned}
\]

\section{Velocity Update}

The method \texttt{apply\_overlap\_momentum()} converts the stored momentum
vector into a velocity change:

\[
    \Delta \mathbf{v}_i
    =
    \frac{\Delta \mathbf{p}_i}{m_i}.
\]

Component-wise:

\[
    v_{x,i}^{\text{new}}
    =
    v_{x,i}
    +
    \frac{\texttt{overlap\_momentum\_x}}{m_i},
\]

\[
    v_{y,i}^{\text{new}}
    =
    v_{y,i}
    +
    \frac{\texttt{overlap\_momentum\_y}}{m_i}.
\]

The position update for the current substep happens in \texttt{Base} before
collision processing. Therefore this velocity change affects the next movement
step.

\section{Current Step Order}

At the current rebuild stage, the relevant order is:

\begin{enumerate}
    \item \texttt{Base} predicts the next position using current velocity.
    \item \texttt{Base.move()} flips the active location buffer.
    \item \texttt{Base.detect\_contacts()} fills each particle's contact list.
    \item \texttt{Dynamics.process\_collisions()} computes overlap momentum.
    \item \texttt{Dynamics.apply\_overlap\_momentum()} updates velocity.
    \item \texttt{Base.record\_step\_report()} stores per-particle values.
\end{enumerate}

\section{Estimating Steps Spent In Collision}

For debugging and stability checks, it is useful to estimate how many simulation
steps a pair of particles will spend in overlap. This estimate helps identify
cases where the relative motion is so fast that contact may last only one or
two sampled steps.

Let the frame rate be \(F\) frames per second and let the simulation use \(N_s\)
substeps per frame. The substep duration is:

\[
    \Delta t_s
    =
    \frac{1}{F N_s}.
\]

For two particles with positions \(\mathbf{x}_i,\mathbf{x}_j\), velocities
\(\mathbf{v}_i,\mathbf{v}_j\), and radii \(r_i,r_j\), define the relative
position and velocity:

\[
    \mathbf{x}_{rel}
    =
    \mathbf{x}_j-\mathbf{x}_i,
    \qquad
    \mathbf{v}_{rel}
    =
    \mathbf{v}_j-\mathbf{v}_i.
\]

The pair is geometrically in contact while:

\[
    \|\mathbf{x}_{rel}+\mathbf{v}_{rel}t\|
    \le
    r_i+r_j.
\]

Expanding this condition gives a quadratic:

\[
    at^2+bt+c \le 0,
\]

where:

\[
\begin{aligned}
    a &= \mathbf{v}_{rel}\cdot\mathbf{v}_{rel},\\
    b &= 2(\mathbf{x}_{rel}\cdot\mathbf{v}_{rel}),\\
    c &= \mathbf{x}_{rel}\cdot\mathbf{x}_{rel}-(r_i+r_j)^2.
\end{aligned}
\]

If the discriminant

\[
    D=b^2-4ac
\]

is negative, the constant-velocity paths do not overlap. If \(D\ge 0\), the
entry and exit times are:

\[
    t_{entry}
    =
    \frac{-b-\sqrt{D}}{2a},
    \qquad
    t_{exit}
    =
    \frac{-b+\sqrt{D}}{2a}.
\]

The continuous collision duration is:

\[
    T_c
    =
    t_{exit}-t_{entry}.
\]

The estimated number of simulation substeps spent in collision is:

\[
    n_c
    =
    T_c F N_s
    =
    \frac{T_c}{\Delta t_s}.
\]

\subsection{Speed-Limit Estimate}

If the only known velocity is a system speed limit \(v_{max}\), the largest
possible relative speed for two particles moving directly toward each other is:

\[
    v_{rel,max}=2v_{max}.
\]

The shortest head-on geometric contact duration is therefore:

\[
    T_{c,min}
    =
    \frac{2(r_i+r_j)}{v_{rel,max}}
    =
    \frac{r_i+r_j}{v_{max}}.
\]

The corresponding conservative lower bound on substeps in collision is:

\[
    n_{c,min}
    =
    \frac{r_i+r_j}{v_{max}} F N_s.
\]

For equal-radius particles \(r_i=r_j=r\):

\[
    n_{c,min}
    =
    \frac{2r}{v_{max}} F N_s.
\]

This estimate ignores collision response. It only answers how many sampled
substeps the particles would overlap if they continued at constant velocity.

\subsection{Sampling Risk}

A useful diagnostic is the relative distance traveled in one substep:

\[
    \Delta x_{rel,s}
    =
    \|\mathbf{v}_{rel}\|\Delta t_s.
\]

If:

\[
    \Delta x_{rel,s} > r_i+r_j,
\]

then the particles move more than the contact radius in one substep. If:

\[
    \Delta x_{rel,s} > 2(r_i+r_j),
\]

then they can cross the entire geometric collision window in one substep. These
cases are risky for overlap-only contact detection because a collision may be
sampled for too few steps, or missed entirely depending on the sampled
positions.

\subsection{Response-Aware Estimate}

The actual number of steps in contact can differ from the geometric estimate
because \texttt{momentum\_per\_area} and
\texttt{inverse\_square\_softening} change particle velocities during overlap.
For each substep, the current model computes:

\[
    J_{ij}
    =
    k_i A_{ij} w(d),
\]

with:

\[
    k_i = \texttt{momentum\_per\_area},
    \qquad
    w(d)=\frac{1}{\max(d^2,s^2)},
    \qquad
    s=\texttt{inverse\_square\_softening}.
\]

The response velocity change is:

\[
    \Delta \mathbf{v}_i
    =
    \frac{-J_{ij}\mathbf{n}_{ij}}{m_i}.
\]

Because this velocity change alters later positions, a response-aware collision
step count is best estimated by running the same substep loop as the simulation:

\begin{enumerate}
    \item compute the current center distance,
    \item detect whether \(d < r_i+r_j\),
    \item compute overlap area \(A_{ij}\),
    \item compute \(J_{ij}=k_iA_{ij}w(d)\),
    \item update velocity from momentum,
    \item move by the substep duration \(\Delta t_s\),
    \item count each substep where overlap is present.
\end{enumerate}

The geometric estimate and the response-aware estimate should both be reported.
The geometric estimate shows whether the time resolution is adequate; the
response-aware estimate shows how the current momentum parameters change the
actual contact duration.

\subsection{Issues}

The are a number of issues that arise while designing a simulation.

\begin{parnumbers}
There is an python menu driven application velocity\_menu.py that takes parameters of the maximum speed and radius of particle which produces the number of frames the particle will be in collision. If there are errors (tunneling etc.) it will report these also.

Reducing the particle radius requires increasing momentum per area and increasing the radius requires lowering it.

\end{parnumbers}



\section{What This Class Does Not Do}

The current \texttt{Dynamics} class does not include bonding, field dynamics,
long-term internal momentum release, classical restitution coefficient \(e\),
or pair-store bookkeeping. It is only the minimal overlap-momentum collision
kernel.

It also does not update both particles inside a single pair operation. That is
intentional for the current particle-owned memory model.

\section{Future Porting Notes}

The most direct C++ or GLSL version would give each particle one worker. That
worker would:

\begin{enumerate}
    \item read particle \(i\)'s contact count and contact array,
    \item loop through the valid contact slots,
    \item read particle \(j\)'s current position and radius,
    \item compute \(A_{ij}\), \(J_{ij}\), and \(-J_{ij}\mathbf{n}_{ij}\),
    \item write only particle \(i\)'s momentum outputs and updated velocity.
\end{enumerate}

The current Python class is shaped to keep that translation straightforward.
\section{Entry Flow}

The field-particle system models inlet flow by activating particles from a preallocated reservoir at a prescribed mass-flow or particle-flow rate. Particles are emitted through a geometric inlet boundary representing the entrance of a pipe or nozzle into the computational domain.

Unlike periodic boundary conditions, particles leaving the domain are not wrapped to the opposite side. Instead, particles that escape the computational limits are returned to an inactive reservoir state and may later be reintroduced through the inlet. This allows the system to model transient accumulation, recirculation, and storage of mass within the domain.

The total system mass is therefore represented by

\[
M_{total}
=
M_{reservoir}
+
M_{stored}
+
M_{out}
\]

where:

\begin{itemize}
    \item $M_{reservoir}$ is the inactive particle mass available for reinjection,
    \item $M_{stored}$ is the mass currently contained within the computational domain,
    \item $M_{out}$ is the cumulative escaped mass.
\end{itemize}

The transient mass balance becomes

\[
\dot{M}_{in}
=
\frac{dM_{stored}}{dt}
+
\dot{M}_{out}
\]

or equivalently,

\[
M_{in}(t)
=
M_{stored}(t)
+
M_{out}(t)
\]

for a domain initially containing no particles.

Particles are introduced at a prescribed particle rate,

\[
\dot{N}_{in}
=
\frac{\dot{M}_{in}}{m_p}
\]

where $m_p$ is the mass represented by an individual particle. For a timestep $\Delta t$, the number of particles introduced is

\[
N_{inject}
=
\dot{N}_{in}\Delta t
\]

Since $N_{inject}$ may not be an integer, an accumulator is used to preserve the correct long-term average injection rate.

Particles are emitted across the pipe cross-section according to

\[
(y-y_c)^2 + (z-z_c)^2 \le R^2
\]

where $(y_c,z_c)$ is the pipe center and $R$ is the pipe radius.

For laminar inflow conditions, the inlet velocity profile may be prescribed using the classical parabolic distribution,

\[
u(r)
=
u_{max}
\left(
1-\frac{r^2}{R^2}
\right)
\]

where

\[
u_{max}=2u_{avg}
\]

and

\[
r^2=(y-y_c)^2+(z-z_c)^2
\]

Alternatively, stochastic perturbations may be added to the inlet velocity to model turbulent or smoke-like inflow behavior.

Inactive reservoir particles immediately terminate execution and therefore do not participate in motion integration, broad-phase traversal, or collision detection. This allows all particles required by the simulation to remain preallocated while only active particles contribute computational cost.

\subsection{Randomized Inlet Distribution}

To avoid artificial lattice alignment and directional bias at the inlet, particles are distributed randomly across the pipe cross-section. A uniformly distributed random point within the circular inlet is generated using polar coordinates.

The radial and angular coordinates are sampled according to

\[
r = R\sqrt{\xi_1}
\]

\[
\theta = 2\pi \xi_2
\]

where $\xi_1,\xi_2 \in [0,1]$ are uniformly distributed random numbers.

The particle position at the inlet plane is then given by

\[
x = x_{in}
\]

\[
y = y_c + r\cos\theta
\]

\[
z = z_c + r\sin\theta
\]

where:

\begin{itemize}
    \item $x_{in}$ is the inlet plane location,
    \item $(y_c,z_c)$ is the center of the pipe cross-section,
    \item $R$ is the pipe radius.
\end{itemize}

The square-root scaling in the radial sampling is necessary to preserve uniform particle density across the circular inlet area. Without this correction, particles would cluster near the pipe center.

The corresponding implementation is

\begin{verbatim}
float r = R * sqrt(rand01());
float theta = 2.0f * PI * rand01();

float y = yc + r * cos(theta);
float z = zc + r * sin(theta);
float x = xin;
\end{verbatim}
\subsection{Development of Entry Flow}

If particles are introduced into the pipe with an initially uniform inlet velocity profile, the no-slip condition imposed at the pipe walls causes the velocity field to evolve naturally toward the classical parabolic fully-developed profile.

At the inlet,

\[
u(r,x=0)=u_0
\]

where $u_0$ is constant across the pipe cross-section.

The pipe wall imposes the no-slip condition,

\[
u(R,x)=0
\]

which progressively slows particles near the wall through particle-wall interactions and local momentum exchange. As the flow advances downstream, a momentum boundary layer develops from the wall toward the pipe centerline.

Initially, the velocity profile remains approximately uniform except near the walls. With increasing downstream distance, the boundary layers grow and eventually merge at the pipe centerline, producing the fully-developed laminar profile,

\[
u(r)
=
u_{max}
\left(
1-\frac{r^2}{R^2}
\right)
\]

This process represents the classical hydrodynamic entrance region of pipe flow.

Within the field-particle framework, the parabolic profile is not imposed analytically. Instead, it emerges naturally from the local particle dynamics, wall interactions, and momentum redistribution within the system.

Consequently, the inlet condition may be prescribed as a simple uniform velocity field while the downstream flow self-organizes into the expected continuum-like behavior.

\end{document}
